\documentclass{article}
\usepackage{graphicx} % Required for inserting images
\usepackage{amsmath}
\usepackage{}

\title{Nombres complexos}
\author{Lluis Gibert i Liam Korsak}
\date{23 de març de 2026}

\begin{document}

\maketitle
\section*{Exercici 2}

\textbf{Demostració:}

Si $z_1$ i $z_2$ són arrels d’un polinomi amb coeficients reals, aleshores
$z_2=\overline{z_1}$.

\[
z_1^n + z_2^n = z_1^n + \overline{z_1}^n
\]

Com que $\overline{z_1^n} = \overline{z_1}^n$, tenim:

\[
z_1^n + z_2^n = z_1^n + \overline{z_1^n}
= 2\operatorname{Re}(z_1^n) \in \mathbb{R}
\]

\textbf{Cas concret:}

\[
x^2 - 2x + 2 = 0
\]

\[
x = 1 \pm i
\]

\[
z_1 = 1+i,\quad z_2 = 1-i
\]

\[
z_1 = \sqrt{2}e^{i\pi/4},\quad z_2 = \sqrt{2}e^{-i\pi/4}
\]

\[
z_1^n = (\sqrt{2})^n e^{i n\pi/4},\quad
z_2^n = (\sqrt{2})^n e^{-i n\pi/4}
\]

\[
z_1^n + z_2^n =
(\sqrt{2})^n \left(e^{i n\pi/4} + e^{-i n\pi/4}\right)
\]

\[
= 2(\sqrt{2})^n \cos\left(\frac{n\pi}{4}\right)
\]

\[
\boxed{
z_1^n + z_2^n = 2(\sqrt{2})^n \cos\left(\frac{n\pi}{4}\right)
}
\]

\end{document}